% Limitations: numbered section BEFORE the Conclusion, 2-3 sentences, 35-60
% words, two items, no mitigation, no promise [most]. Items: preferably
% "On the one hand ... On the other hand ..." [current], otherwise
% "First ... In addition ..." [older]. Present tense.
\section{Limitations} \label{sec:limitations}
This work has two main limitations.
On the one hand, <extra cost or stability issue>.
On the other hand, <scope limit>, so whether X transfers to <other domain> remains untested.

% Conclusion: ONE paragraph, 80-140 words, mirrors the abstract, no new claims;
% 3-6 sentences; past tense [most] or present [some] -- one tense per paragraph.
\section{Conclusion} \label{sec:conclusion}
Prior <paradigm> methods such as <Prior> \citep{Placeholder-A} bring <benefit>, but they suffer from <problem 1> caused by <cause> and from <problem 2> because they lack <mechanism>.
We proposed X, which <achieves objective A> through <mechanism A> and <achieves objective B> through <mechanism B> (\ie <core concept A> and <core concept B>).
Experiments on several models confirmed that X is effective and makes <task> more <epithet>.
Future work will extend X to <domain>.
% Optional footnote on the last sentence pointing to a statement (2026-style ICLR kits:
% "LLM use is described in Appendix~\ref{apdx:use-of-llms}"; NeurIPS: broader impacts). ICLR 2027
% discloses LLM use in the AI Use Statement before the references instead.

\label{main-text-end}  % page check: this label's page must equal the page limit (see SKILL.md)

% Camera-ready only (ICLR kit: unnumbered third-level heading; NeurIPS: \begin{ack}):
% \subsubsection*{Acknowledgments}
% This research is supported by <grant>.
% \clearpage
